\subsection{Prime choices passing the compact elliptic tests}
\label{sec:ghost-elliptic-survival}

\begin{proposition}[An original-factorial family passing every point-order test]
\label{prop:ghost-elliptic-passing}
The original A.44 recursion admits choices at every stage with
\[
 3\notin(\mathbb F_{q_j}^{\times})^3,
       \qquad 2\notin(\mathbb F_{r_j}^{\times})^3.
                                                        \tag{21}
\]
Its resulting ordinary infinite-support pairs satisfy every
point-order condition of
Proposition~\ref{prop:ghost-elliptic-support-test}.
This statement constructs the pairs and proves these tests;
it does not assert a section above any of the pairs.
\end{proposition}
\begin{proof}
We first prove the field intersection needed to retain all
of A.44's original congruences. For a cyclotomic field
\(F\) containing \(\mu_3\), the map
\[
 \mathbb Q^\times/(\mathbb Q^\times)^3
       \longrightarrow F^\times/(F^\times)^3
                                                        \tag{22}
\]
is injective. A cube root of a rational noncube would put
its real pure cubic field in the abelian extension
\(F/\mathbb Q\). That degree-three real field is not
normal, since its other two embeddings are nonreal;
every intermediate field of an abelian extension is
normal. This proves (22), including either sign of the
rational number.

At the \(q\)-choice of a stage, retain the notation
\(a,b,n,m,M,Q_0\) of A.44, and take
\(F=\mathbb Q(\mu_M)\); here \(n\mid m\mid M\).
The original selection gives an automorphism of
\(R=F(b^{1/n})\) restricting to the prescribed
\(Q_0\)-automorphism on \(F\) and fixing
\(\mathbb Q(\mu_n,b^{1/n})\).
The classes of \(3\) and \(b\) are independent in
\(\mathbb Q^\times/(\mathbb Q^\times)^3\): the
valuations are \(v_3(b)=0\), \(v_2(b)=1\).
They remain independent by (22). Since \(R/F\)
is cyclic and \(b\) is not a cube in \(F\), its
unique cubic subextension is \(F(b^{1/3})\).
It follows that
\[
                 R\cap F(3^{1/3})=F.                 \tag{23}
\]
The prescribed automorphism of \(F\) fixes \(\mu_3\),
because \(Q_0\equiv1\pmod m\). Thus on
\(F(3^{1/3})\) it has extensions sending
\(3^{1/3}\) to \(\zeta_3^e3^{1/3}\),
for each \(e\in\mathbb F_3\). Equation (23)
allows either nonzero choice to be combined with the
original automorphism of \(R\). Chebotarev
\cite[Theorem 8.31]{GhostMilneANT} gives infinitely
many primes in its conjugacy class. They retain the
full CRT congruence and original radical splitting,
and make \(3\) a noncube. The noncube condition is
invariant under conjugation, which only multiplies
\(e\) by a unit. We can therefore retain every
lower bound and every exclusion of old support primes.

At the subsequent \(r\)-choice put
\(F=\mathbb Q(\mu_{Mq})\) and \(c=aq^n\).
Now \(v_2(c)=0\), \(v_3(c)=1\), so the cube
classes of \(2,c\) are independent, and
\[
        F(c^{1/m})\cap F(2^{1/3})=F.
\]
The same argument combines the old CRT and radical
automorphism with a nontrivial action on \(2^{1/3}\).
It selects \(r\) with \(2\) a noncube. Induction
therefore carries out A.44 with (21) at every stage.
All of its ordinary local convergence and infinite-support
properties are retained, since no original condition has
been discarded.

It remains to prove the stated finite elliptic tests.
Every new prime is \(1\bmod3\). The old radical
splitting gives \(-2\) a cube at \(q_j\), since
\(b=-2\) times a rational cube and \(3\mid n\).
It gives \(3\) a cube at \(r_j\), since
\(aq^n=3\) times a rational cube and \(3\mid m\).
For \(E':y^2=x^3+9\), the nonconstant part of
the three-division polynomial is \(x^3+36\).
Over \(\mathbb F_p\) with \(p\equiv1\pmod3\),
the corresponding ordinates \(\pm3\sqrt{-3}\)
are present. Thus its full three-torsion is rational
exactly when \(-36\) is a cube. At either kind of
prime in (21), one of \(2,3\) is a cube and the
other is not, so \(-36\) is not a cube. Consequently
\[
                 E'(\mathbb F_p)[3]=\langle T\rangle.
                                                        \tag{24}
\]

In fact the entire three-primary subgroup in (24)
has order three. It is cyclic because its subgroup
killed by three is cyclic. Let
\(\rho(x,y)=(\zeta_3x,y)\), an automorphism over
\(\mathbb F_p\). The horizontal chord through
\((x,y),(\zeta_3x,y),(\zeta_3^2x,y)\) proves
\(1+\rho+\rho^2=0\) as an elliptic endomorphism.
On a cyclic group of order at least nine, \(\rho\)
would therefore give a root of \(X^2+X+1\bmod9\).
There is no such root: a root modulo three must be
\(1\), and substitution of \(1+3k\) gives
\(3\bmod9\). This proves
\(E'(\mathbb F_p)\{3\}=\langle T\rangle\).

The curves \(E,E'\) have the same number of points
over this finite field, as follows from their isogeny
\cite[V.4.1]{GhostMilneEC}. Since the kernel of
\(\psi\) is \(\mu_3\), all three of its points
are rational. Its image \(\psi E(\mathbb F_p)\)
therefore has index three and order prime to three.
The exact Kummer formula of
Proposition~\ref{prop:ghost-elliptic-selmer} gives
\[
 \delta_\psi(P_0+T)=-2,
        \qquad \delta_\psi(P_0-T)=-9.
\]
The first value is a cube at every \(q_j\), and
the second is a cube at every \(r_j\). The indicated
points belong to \(\psi E(\mathbb F_p)\) and hence
have prime-to-three order. These are precisely all
the tests in (16).
\end{proof}

As small independent certificates, \(q=2281\) has
\(-2\) a twenty-fourth power and \(3\) a noncube;
\(\#E'(\mathbb F_q)=2361\) and
\(\operatorname{ord}((-2,1))=787\).
The prime \(r=19441\) has \(3\) a one-hundred-and-twentieth
power and \(2\) a noncube;
\(\#E'(\mathbb F_r)=19200\) and
\(\operatorname{ord}((3,-6))=40\).
These values are checked in
\texttt{check\_elliptic\_descent.py}. They illustrate the
finite-field conditions; they omit the enormous recursive
CRT moduli and are not asserted to be the actual A.44 primes.

The precise unresolved case of the compact elliptic route
is therefore nonempty: it contains the constructed family
of this proposition. No full section above this family
has been constructed, and no proof here excludes it by
the infinite fundamental group of the open elliptic cover.
